\section{Rectification}
\subsection{Identification des appels récursifs}

L'identification des appels récursifs est assez complexe lorsqu'on considère des
fonctions mutuellement récursives ou des fonctions qui prennent d'autres
fonctions en argument. On introduit alors la notion
d'\textit{ensemble rectifiant :}

\begin{definition}
    Un ensemble de noms de variables $E$ \textit{rectifie} une forme linéaire
    $l$ lorsque tout appel de fonction potentiellement récursif dans $l$
    (directement ou indirectement) est un appel à une fonction dont le nom est
    dans $E$ (l'opération élémentaire <<~appel de fonction~>> s'applique sur un
    nom de variable).
\end{definition}

\begin{example} \strut
    La forme linéaire de la partie précédente correspondant au corps de la
    fonction \texttt{sum} est rectifiable par $\{\texttt{sum}, a_6\}$. On peut
    remarquer qu'il est aussi rectifiable par
    $\{\texttt{sum}, a_6, \texttt{l}, a_4, a_9\}$; l'ajout de nouveaux noms de
    variables n'altère pas le caractère rectifiant de l'ensemble.
\end{example}

On cherche à avoir des ensembles rectifiants les plus petits possibles, car
nous ne pouvons pas modifier le code de certaines fonctions (de la bibliothèque
standard \texttt{OCaml}, par exemple). On peut alors considérer le problème de
décision suivant :

\begin{definition}
    Le problème de décision ENS\_RECTIFIANT consiste à déterminer, étant donné
    une forme linéaire $l$ et un entier $k$, s'il existe un ensemble rectifiant
    $l$ de cardinal plus petit que $k$.
\end{definition}

\begin{proposition}
    Le problème ENS\_RECTIFIANT est indécidable.
\end{proposition}

\begin{proof}
    Par réduction depuis le co-problème de l'arrêt.
\end{proof}

Notre algorithme pour calculer un ensemble rectifiant minimal pour une forme
linéaire n'est alors qu'une approximation, et il existe des entrées sur lesquels
il abandonne. Cependant il fonctionne bien sur la plupart des programmes
rencontrés en pratique. Il est implémenté par saturation, avec des règles telles
que <<~si $f \in E$ et s'il existe une instruction
$a_k := f$, alors $a_k \in E$~>>, avec un ensemble initialement réduit aux noms
des fonctions à rectifier.

\subsection{Algorithme de rectification}

Avec les appels récursifs ainsi identifiés, on a alors tous les éléments pour
mettre en place l'algorithme décrit dans la section 1.

\begin{algorithm}[H]
    \caption{Rectifier}
    \Entree{
        Une forme linéaire $l$ ne contenant pas de variables nommées
        \texttt{cont} ou \texttt{new\_cont} et $E$ un ensemble rectifiant $l$.
    }
    \Sortie{
        Une forme linéaire récursive terminale $l'$ sous CPS équivalente à $l$
    }

    $(a, e) \leftarrow$ le premier élément de $l$ \\
    $q \leftarrow$ le reste de $l$ \\

    \uSi{$e$ contient une définition d'une fonction de $E$}{
        Mettre à jour $e$ tel que la fonction définie dedans:
        \begin{itemize}
            \item Prenne un argument \texttt{cont} en plus
            \item Évalue $Rectifier(l', E)$ lors de son appel o\`u $l'$
                  est la forme linéaire évaluée \\ avant cette redéfinition.
        \end{itemize}
    }

    \uSi{$e$ ne contient pas d'appel à une fonction de $E$}{
        \uSi{$q$ est vide}{
            Renvoyer $\left\{ \hspace{0.2cm} \begin{aligned}
                    a   & := e                & \\
                    res & := \texttt{cont } a &
                \end{aligned} \right. $
        }\uSinon{
            $q_{rec} \leftarrow$ Rectifier($q$, $E$) \\
            Renvoyer $(a, e) \cdot q_{rec}$
        }
    }
    \uSinon{
        \uSi{$e$ est un appel de fonction}{
            $e_{rec} \leftarrow$ $e$ avec \texttt{cont} rajouté en dernier
            argument
        }\uSinon{
            $e_{rec} \leftarrow$ $e$ avec chaque enfant localement terminal $c$
            remplacé par Rectifier($c$, $E$)
        }

        \bigskip

        \uSi{$q$ est vide}{
            Renvoyer $[(a, e_{rec})]$
        }\uSinon{
            $q_{rec} \leftarrow$ Rectifier($q$, $E$) \\
            Renvoyer $\left\{ \hspace{0.2cm} \begin{aligned}
                    \texttt{new\_cont} & := \texttt{fun } a \texttt{ -> } q_{rec} & \\
                    \texttt{cont}      & := \texttt{new\_cont}                    & \\
                    a                  & := e_{rec}                               &
                \end{aligned} \right.$
        }
    }
\end{algorithm}

On peut finalement démontrer la correction de cet algorithme. Le caractère
récursif terminal de la forme linéaire renvoyée se déduit immédiatement de
l'algorithme en constatant que les appels récursifs sont toujours à la fin d'une
forme linéaire, et l'équivalence entre $l$ et $l'$ peut se démontrer par
récurrence sur le nombre d'opérations élémentaires effectuées lors d'une
exécution particulière de $l$, en montrant que $l'$ a alors le même résultat (et
les mêmes effets de bord).

\subsection{Utilisation d'accumulateurs}

Une dernière amélioration à faire est de constater que dans certains cas, une
fonction sous CPS peut être remplacée par une fonction utilisant un accumulateur
pour une exécution plus efficace en mémoire. C'est notamment le cas si :
\begin{itemize}
    \item Toute nouvelle continuation dans la fonction est la composition de la
          continuation précédente et d'une fonction $f_i$~;
    \item La fonction originale a une seule valeur de retour $init$ en dehors
          des appels récursifs~;
    \item Les fonctions $f_i$ commutent.
\end{itemize}

Dans ce cas, la continuation du premier appel de la fonction est forcément
appelé avec pour argument $f_1(f_2(\ldots f_n(init)))$ qui est égal à
$f_n(\ldots f_2(f_1(init)))$ puisque les fonctions $f_i$ commutent, et on
remarque que $f_1(init)$ peut être calculé dès le premier appel récursif car
$f_1$ et $init$ sont tous deux connus dès ce premier appel. On peut alors
remplacer la continuation par un accumulateur contenant la valeur de
$f_k(f_{k-1}(\ldots f_1(init)))$ au $k$-ième appel récursif, amélioration que
j'ai aussi implémenté.
